% =============================================================================
% investigacion/p6-open-layers/notas.tex — INV-P6 (issue #12)
% Integrante: Enrique A Gonzalez Moreira
% Fragmento puro: se incluye desde preguntas/p6-ncapas-calidad/solucion.tex
% mediante % =============================================================================
% investigacion/p6-open-layers/notas.tex — INV-P6 (issue #12)
% Integrante: Enrique A Gonzalez Moreira
% Fragmento puro: se incluye desde preguntas/p6-ncapas-calidad/solucion.tex
% mediante % =============================================================================
% investigacion/p6-open-layers/notas.tex — INV-P6 (issue #12)
% Integrante: Enrique A Gonzalez Moreira
% Fragmento puro: se incluye desde preguntas/p6-ncapas-calidad/solucion.tex
% mediante % =============================================================================
% investigacion/p6-open-layers/notas.tex — INV-P6 (issue #12)
% Integrante: Enrique A Gonzalez Moreira
% Fragmento puro: se incluye desde preguntas/p6-ncapas-calidad/solucion.tex
% mediante \input{investigacion/p6-open-layers/notas}.
% Estructura del issue #12: \section*{Notas de investigación: …} + Síntesis +
% Conceptos clave (tabla booktabs) + Puntos de conexión + Notas de diseño.
% Fuente de trabajo: notas.md (misma carpeta).
% Citas: páginas impresas de Richards (2015) y Martin (2017); Richards y Ford
% (2020) no conserva paginación al extraer el texto del PDF y se cita por
% capítulo; Cockburn (2005) y Palermo (2008) son web sin paginación.
% Todas las claves están en core.bib (richards2015, richardsford2020,
% martin2017, cockburn2005, palermo2008): fuentes.bib no declara nada.
% Comillas: se usan comillas tipográficas literales, no \enquote, para no
% depender de csquotes (no está en la lista de paquetes de preamble.tex, #13).
% =============================================================================

% Grupo local: \emergencystretch evita líneas desbordadas sin afectar a otros capítulos.
\begingroup\emergencystretch=3em \hbadness=2000
\section*{Notas de investigación: Open Layers, Sinkhole Anti-pattern y atributos de calidad en N-Capas}

\section{Síntesis}

La arquitectura N-Capas organiza el sistema en capas horizontales en las que cada capa conoce
únicamente a la inmediatamente inferior. Una capa es \textbf{cerrada} cuando ninguna solicitud
puede saltársela —toda petición debe atravesarla para alcanzar la siguiente— y \textbf{abierta}
cuando sí puede evitarse, accediendo directamente a una capa más profunda
\parencite[3]{richards2015}. Las capas cerradas son las que materializan el \emph{aislamiento de
capas}: los cambios realizados en una capa no impactan a los componentes de las demás mientras
se respeten los contratos, de modo que cada capa es independiente y conoce poco o nada del
interior de las otras \parencite[3-4]{richards2015}.

Abrir una capa compra rendimiento a costa de ese aislamiento. El caso típico es el
\emph{fast-lane reader}: se declaran abiertas las capas de negocio y persistencia para que la
presentación lea directamente la base de datos en consultas simples de recuperación
\parencite[cap.~10]{richardsford2020}. El precio es un acoplamiento transversal —un cambio en
persistencia pasa a impactar a la vez a negocio y a presentación— y una arquitectura
descrita como frágil y cara de cambiar \parencite[4]{richards2015}. Por eso la decisión es
sobre todo de gobernanza: no documentar ni comunicar qué capas son abiertas y cuáles cerradas,
y por qué, produce arquitecturas acopladas y frágiles, difíciles de probar, mantener y desplegar
\parencite[5]{richards2015}.

El \emph{architecture sinkhole anti-pattern} aparece cuando las solicitudes atraviesan las capas
como simple \emph{pass-through}, sin que ninguna aporte lógica de negocio: el coste es
instanciación y procesamiento innecesarios, con impacto en memoria y rendimiento
\parencite[cap.~10]{richardsford2020}. Importa la dirección causal: las fuentes no presentan las
capas abiertas como causa del sinkhole, sino como reacción propuesta cuando el sinkhole domina
el tráfico, con el aviso explícito de que entonces será más difícil controlar el cambio por la
falta de aislamiento \parencite[7]{richards2015}. La \textbf{regla del 80/20} es el criterio de
diagnóstico: es normal que en torno al 20\,\% de las peticiones sean \emph{pass-through} y el
80\,\% porten lógica, y si la proporción se invierte eso indica que la arquitectura en capas no
es el estilo correcto para ese dominio \parencite[7]{richards2015}.

Las calificaciones bajas de la N-Capas en agilidad, facilidad de despliegue y escalabilidad no
son accidentales: proceden de su naturaleza monolítica y del acoplamiento entre componentes
\parencite[8-9]{richards2015}. Para requisitos de alta escalabilidad y agilidad, la Hexagonal y
la Limpia son estructuralmente más adecuadas porque invierten la dirección de las dependencias:
el núcleo declara los puertos y la infraestructura los implementa \parencite{cockburn2005}, con
dependencias que apuntan solo hacia adentro \parencite[203]{martin2017} y un dominio que no se
acopla más que a sí mismo \parencite{palermo2008}.

\section{Conceptos clave}

\begin{table}[htbp]
  \centering
  \footnotesize
  \renewcommand{\arraystretch}{1.2}
  \begin{tabular}{@{}p{0.22\linewidth}p{0.46\linewidth}p{0.24\linewidth}@{}}
    \toprule
    \textbf{Concepto} & \textbf{Definición} & \textbf{Fuente} \\
    \midrule
    \raggedright Capa cerrada
      & \raggedright Capa que no puede saltarse: toda solicitud debe atravesarla para llegar a
        la capa inferior
      & \raggedright \textcite[3]{richards2015}; \textcite[cap.~10]{richardsford2020}
        \tabularnewline
    \raggedright Capa abierta
      & \raggedright Capa que una solicitud puede evitar, accediendo directamente a una capa más
        profunda
      & \raggedright \textcite[3]{richards2015}; \textcite[cap.~10]{richardsford2020}
        \tabularnewline
    \raggedright Aislamiento de capas
      & \raggedright Los cambios en una capa no impactan a las demás mientras se mantengan los
        contratos entre capas
      & \raggedright \textcite[3-4]{richards2015}; \textcite[cap.~10]{richardsford2020}
        \tabularnewline
    \raggedright Gobernanza de las capas abiertas
      & \raggedright No documentar qué capas son abiertas o cerradas produce arquitecturas
        acopladas y frágiles, difíciles de probar, mantener y desplegar
      & \raggedright \textcite[5]{richards2015} \tabularnewline
    \raggedright \emph{Fast-lane reader}
      & \raggedright Declarar abiertas negocio y persistencia para que la presentación lea
        directamente la base de datos en consultas simples de alto volumen
      & \raggedright \textcite[cap.~10]{richardsford2020} \tabularnewline
    \raggedright \emph{Architecture sinkhole anti-pattern}
      & \raggedright Solicitudes que atraviesan las capas como \emph{pass-through}, sin lógica
        en ninguna, con coste de instanciación y procesamiento
      & \raggedright \textcite[7]{richards2015}; \textcite[cap.~10]{richardsford2020}
        \tabularnewline
    \raggedright Regla del 80/20
      & \raggedright Si en torno al 80\,\% de las peticiones son \emph{pass-through}, el estilo
        en capas no es el adecuado para ese dominio
      & \raggedright \textcite[7]{richards2015}; \textcite[cap.~10]{richardsford2020}
        \tabularnewline
    \raggedright Partición técnica frente a partición por dominio
      & \raggedright En N-Capas el dominio queda disperso entre capas técnicas, por lo que un
        enfoque dirigido por el dominio encaja mal
      & \raggedright \textcite[cap.~10]{richardsford2020} \tabularnewline
    \raggedright Calificaciones de la N-Capas
      & \raggedright Agilidad y facilidad de despliegue bajas por la naturaleza monolítica y el
        acoplamiento; escalabilidad baja por la misma razón
      & \raggedright \textcite[8-9]{richards2015}; \textcite[cap.~10]{richardsford2020}
        \tabularnewline
    \raggedright Regla de Dependencia
      & \raggedright Las dependencias del código fuente apuntan solo hacia adentro, hacia las
        políticas de mayor nivel
      & \raggedright \textcite[203]{martin2017} \tabularnewline
    \raggedright Testabilidad del núcleo
      & \raggedright Las reglas de negocio se prueban sin interfaz de usuario, base de datos ni
        servidor
      & \raggedright \textcite[202]{martin2017} \tabularnewline
    \raggedright Puertos y adaptadores; dominio al centro
      & \raggedright El núcleo declara los puertos y la infraestructura los implementa; el
        Modelo de Dominio no se acopla más que a sí mismo
      & \raggedright \textcite{cockburn2005}; \textcite{palermo2008} \tabularnewline
    \bottomrule
  \end{tabular}
  \caption{Conceptos clave de la investigación INV-P6.}
  \label{tab:inv-p6-conceptos}
\end{table}

\section{Puntos de conexión con las preguntas}

\begin{itemize}
  \item \textbf{P6, primer apartado:} capa cerrada frente a capa abierta y aislamiento de capas
        dan el mecanismo exacto por el que abrir una capa degrada la mantenibilidad
        \parencite[3-5]{richards2015}.
  \item \textbf{P6, segundo apartado:} el sinkhole y la regla del 80/20 permiten mostrar que la
        apertura de capas es reacción al antipatrón, no su causa, y que la regla es un
        diagnóstico de encaje del estilo, no una perilla de optimización
        \parencite[7]{richards2015}.
  \item \textbf{P6, tercer apartado:} las calificaciones bajas \parencite[8-9]{richards2015} y
        la partición técnica \parencite[cap.~10]{richardsford2020} se contraponen a la inversión
        de dependencias de la Hexagonal y la Limpia \parencite[203]{martin2017}.
  \item \textbf{P1 y P2:} el mismo contraste N-Capas frente a Hexagonal/Onion sostiene el
        análisis de \emph{trade-offs} de P1 y la comparación de simetría de P2.
\end{itemize}

\section*{Notas de diseño}

\begin{itemize}
  \item \textbf{Dirección causal del sinkhole.} El enunciado de P6 sugiere que las capas
        abiertas llevan al sinkhole; las fuentes dicen lo inverso: son la reacción propuesta
        cuando el sinkhole domina \parencite[7]{richards2015}. La investigación documenta el
        bucle completo —capas vacías, sinkhole, apertura sancionada del rodeo, vaciado
        acelerado— para que la respuesta pueda ser fiel a las fuentes sin contradecir el
        enunciado.
  \item \textbf{Paginación de Richards y Ford (2020).} El PDF de \texttt{Bibliografias/} no
        conserva la paginación impresa al extraer el texto, así que ese libro se cita por
        capítulo y sección. Las páginas de \textcite{richards2015} se verificaron contra la
        numeración impresa del capítulo~1, que en el PDF local aparece desplazada por las
        páginas preliminares.
  \item \textbf{Doble cita de Richards.} \textcite{richards2015} aporta las calificaciones del
        patrón y \textcite{richardsford2020} el tratamiento ampliado del mismo autor principal
        (capas abiertas, \emph{fast-lane reader}, partición técnica). Se citan juntos donde
        coinciden y por separado donde cada uno es más explícito.
  \item \textbf{De la Torre Llorente y otros (2010).} Figura en \texttt{fuentes.md} como
        contexto del seminario sobre N-Capas orientada al dominio, pero no aporta ninguna
        afirmación de estas notas, así que no se cita.
  \item \textbf{DDD sin Evans.} El enunciado menciona principios de DDD. Se sostienen con la
        partición por dominio \parencite[cap.~10]{richardsford2020} y el dominio al centro
        \parencite{palermo2008}, sin recurrir a Evans, que no figura en la bibliografía
        consultada para esta investigación.
\end{itemize}
\endgroup
.
% Estructura del issue #12: \section*{Notas de investigación: …} + Síntesis +
% Conceptos clave (tabla booktabs) + Puntos de conexión + Notas de diseño.
% Fuente de trabajo: notas.md (misma carpeta).
% Citas: páginas impresas de Richards (2015) y Martin (2017); Richards y Ford
% (2020) no conserva paginación al extraer el texto del PDF y se cita por
% capítulo; Cockburn (2005) y Palermo (2008) son web sin paginación.
% Todas las claves están en core.bib (richards2015, richardsford2020,
% martin2017, cockburn2005, palermo2008): fuentes.bib no declara nada.
% Comillas: se usan comillas tipográficas literales, no \enquote, para no
% depender de csquotes (no está en la lista de paquetes de preamble.tex, #13).
% =============================================================================

% Grupo local: \emergencystretch evita líneas desbordadas sin afectar a otros capítulos.
\begingroup\emergencystretch=3em \hbadness=2000
\section*{Notas de investigación: Open Layers, Sinkhole Anti-pattern y atributos de calidad en N-Capas}

\section{Síntesis}

La arquitectura N-Capas organiza el sistema en capas horizontales en las que cada capa conoce
únicamente a la inmediatamente inferior. Una capa es \textbf{cerrada} cuando ninguna solicitud
puede saltársela —toda petición debe atravesarla para alcanzar la siguiente— y \textbf{abierta}
cuando sí puede evitarse, accediendo directamente a una capa más profunda
\parencite[3]{richards2015}. Las capas cerradas son las que materializan el \emph{aislamiento de
capas}: los cambios realizados en una capa no impactan a los componentes de las demás mientras
se respeten los contratos, de modo que cada capa es independiente y conoce poco o nada del
interior de las otras \parencite[3-4]{richards2015}.

Abrir una capa compra rendimiento a costa de ese aislamiento. El caso típico es el
\emph{fast-lane reader}: se declaran abiertas las capas de negocio y persistencia para que la
presentación lea directamente la base de datos en consultas simples de recuperación
\parencite[cap.~10]{richardsford2020}. El precio es un acoplamiento transversal —un cambio en
persistencia pasa a impactar a la vez a negocio y a presentación— y una arquitectura
descrita como frágil y cara de cambiar \parencite[4]{richards2015}. Por eso la decisión es
sobre todo de gobernanza: no documentar ni comunicar qué capas son abiertas y cuáles cerradas,
y por qué, produce arquitecturas acopladas y frágiles, difíciles de probar, mantener y desplegar
\parencite[5]{richards2015}.

El \emph{architecture sinkhole anti-pattern} aparece cuando las solicitudes atraviesan las capas
como simple \emph{pass-through}, sin que ninguna aporte lógica de negocio: el coste es
instanciación y procesamiento innecesarios, con impacto en memoria y rendimiento
\parencite[cap.~10]{richardsford2020}. Importa la dirección causal: las fuentes no presentan las
capas abiertas como causa del sinkhole, sino como reacción propuesta cuando el sinkhole domina
el tráfico, con el aviso explícito de que entonces será más difícil controlar el cambio por la
falta de aislamiento \parencite[7]{richards2015}. La \textbf{regla del 80/20} es el criterio de
diagnóstico: es normal que en torno al 20\,\% de las peticiones sean \emph{pass-through} y el
80\,\% porten lógica, y si la proporción se invierte eso indica que la arquitectura en capas no
es el estilo correcto para ese dominio \parencite[7]{richards2015}.

Las calificaciones bajas de la N-Capas en agilidad, facilidad de despliegue y escalabilidad no
son accidentales: proceden de su naturaleza monolítica y del acoplamiento entre componentes
\parencite[8-9]{richards2015}. Para requisitos de alta escalabilidad y agilidad, la Hexagonal y
la Limpia son estructuralmente más adecuadas porque invierten la dirección de las dependencias:
el núcleo declara los puertos y la infraestructura los implementa \parencite{cockburn2005}, con
dependencias que apuntan solo hacia adentro \parencite[203]{martin2017} y un dominio que no se
acopla más que a sí mismo \parencite{palermo2008}.

\section{Conceptos clave}

\begin{table}[htbp]
  \centering
  \footnotesize
  \renewcommand{\arraystretch}{1.2}
  \begin{tabular}{@{}p{0.22\linewidth}p{0.46\linewidth}p{0.24\linewidth}@{}}
    \toprule
    \textbf{Concepto} & \textbf{Definición} & \textbf{Fuente} \\
    \midrule
    \raggedright Capa cerrada
      & \raggedright Capa que no puede saltarse: toda solicitud debe atravesarla para llegar a
        la capa inferior
      & \raggedright \textcite[3]{richards2015}; \textcite[cap.~10]{richardsford2020}
        \tabularnewline
    \raggedright Capa abierta
      & \raggedright Capa que una solicitud puede evitar, accediendo directamente a una capa más
        profunda
      & \raggedright \textcite[3]{richards2015}; \textcite[cap.~10]{richardsford2020}
        \tabularnewline
    \raggedright Aislamiento de capas
      & \raggedright Los cambios en una capa no impactan a las demás mientras se mantengan los
        contratos entre capas
      & \raggedright \textcite[3-4]{richards2015}; \textcite[cap.~10]{richardsford2020}
        \tabularnewline
    \raggedright Gobernanza de las capas abiertas
      & \raggedright No documentar qué capas son abiertas o cerradas produce arquitecturas
        acopladas y frágiles, difíciles de probar, mantener y desplegar
      & \raggedright \textcite[5]{richards2015} \tabularnewline
    \raggedright \emph{Fast-lane reader}
      & \raggedright Declarar abiertas negocio y persistencia para que la presentación lea
        directamente la base de datos en consultas simples de alto volumen
      & \raggedright \textcite[cap.~10]{richardsford2020} \tabularnewline
    \raggedright \emph{Architecture sinkhole anti-pattern}
      & \raggedright Solicitudes que atraviesan las capas como \emph{pass-through}, sin lógica
        en ninguna, con coste de instanciación y procesamiento
      & \raggedright \textcite[7]{richards2015}; \textcite[cap.~10]{richardsford2020}
        \tabularnewline
    \raggedright Regla del 80/20
      & \raggedright Si en torno al 80\,\% de las peticiones son \emph{pass-through}, el estilo
        en capas no es el adecuado para ese dominio
      & \raggedright \textcite[7]{richards2015}; \textcite[cap.~10]{richardsford2020}
        \tabularnewline
    \raggedright Partición técnica frente a partición por dominio
      & \raggedright En N-Capas el dominio queda disperso entre capas técnicas, por lo que un
        enfoque dirigido por el dominio encaja mal
      & \raggedright \textcite[cap.~10]{richardsford2020} \tabularnewline
    \raggedright Calificaciones de la N-Capas
      & \raggedright Agilidad y facilidad de despliegue bajas por la naturaleza monolítica y el
        acoplamiento; escalabilidad baja por la misma razón
      & \raggedright \textcite[8-9]{richards2015}; \textcite[cap.~10]{richardsford2020}
        \tabularnewline
    \raggedright Regla de Dependencia
      & \raggedright Las dependencias del código fuente apuntan solo hacia adentro, hacia las
        políticas de mayor nivel
      & \raggedright \textcite[203]{martin2017} \tabularnewline
    \raggedright Testabilidad del núcleo
      & \raggedright Las reglas de negocio se prueban sin interfaz de usuario, base de datos ni
        servidor
      & \raggedright \textcite[202]{martin2017} \tabularnewline
    \raggedright Puertos y adaptadores; dominio al centro
      & \raggedright El núcleo declara los puertos y la infraestructura los implementa; el
        Modelo de Dominio no se acopla más que a sí mismo
      & \raggedright \textcite{cockburn2005}; \textcite{palermo2008} \tabularnewline
    \bottomrule
  \end{tabular}
  \caption{Conceptos clave de la investigación INV-P6.}
  \label{tab:inv-p6-conceptos}
\end{table}

\section{Puntos de conexión con las preguntas}

\begin{itemize}
  \item \textbf{P6, primer apartado:} capa cerrada frente a capa abierta y aislamiento de capas
        dan el mecanismo exacto por el que abrir una capa degrada la mantenibilidad
        \parencite[3-5]{richards2015}.
  \item \textbf{P6, segundo apartado:} el sinkhole y la regla del 80/20 permiten mostrar que la
        apertura de capas es reacción al antipatrón, no su causa, y que la regla es un
        diagnóstico de encaje del estilo, no una perilla de optimización
        \parencite[7]{richards2015}.
  \item \textbf{P6, tercer apartado:} las calificaciones bajas \parencite[8-9]{richards2015} y
        la partición técnica \parencite[cap.~10]{richardsford2020} se contraponen a la inversión
        de dependencias de la Hexagonal y la Limpia \parencite[203]{martin2017}.
  \item \textbf{P1 y P2:} el mismo contraste N-Capas frente a Hexagonal/Onion sostiene el
        análisis de \emph{trade-offs} de P1 y la comparación de simetría de P2.
\end{itemize}

\section*{Notas de diseño}

\begin{itemize}
  \item \textbf{Dirección causal del sinkhole.} El enunciado de P6 sugiere que las capas
        abiertas llevan al sinkhole; las fuentes dicen lo inverso: son la reacción propuesta
        cuando el sinkhole domina \parencite[7]{richards2015}. La investigación documenta el
        bucle completo —capas vacías, sinkhole, apertura sancionada del rodeo, vaciado
        acelerado— para que la respuesta pueda ser fiel a las fuentes sin contradecir el
        enunciado.
  \item \textbf{Paginación de Richards y Ford (2020).} El PDF de \texttt{Bibliografias/} no
        conserva la paginación impresa al extraer el texto, así que ese libro se cita por
        capítulo y sección. Las páginas de \textcite{richards2015} se verificaron contra la
        numeración impresa del capítulo~1, que en el PDF local aparece desplazada por las
        páginas preliminares.
  \item \textbf{Doble cita de Richards.} \textcite{richards2015} aporta las calificaciones del
        patrón y \textcite{richardsford2020} el tratamiento ampliado del mismo autor principal
        (capas abiertas, \emph{fast-lane reader}, partición técnica). Se citan juntos donde
        coinciden y por separado donde cada uno es más explícito.
  \item \textbf{De la Torre Llorente y otros (2010).} Figura en \texttt{fuentes.md} como
        contexto del seminario sobre N-Capas orientada al dominio, pero no aporta ninguna
        afirmación de estas notas, así que no se cita.
  \item \textbf{DDD sin Evans.} El enunciado menciona principios de DDD. Se sostienen con la
        partición por dominio \parencite[cap.~10]{richardsford2020} y el dominio al centro
        \parencite{palermo2008}, sin recurrir a Evans, que no figura en la bibliografía
        consultada para esta investigación.
\end{itemize}
\endgroup
.
% Estructura del issue #12: \section*{Notas de investigación: …} + Síntesis +
% Conceptos clave (tabla booktabs) + Puntos de conexión + Notas de diseño.
% Fuente de trabajo: notas.md (misma carpeta).
% Citas: páginas impresas de Richards (2015) y Martin (2017); Richards y Ford
% (2020) no conserva paginación al extraer el texto del PDF y se cita por
% capítulo; Cockburn (2005) y Palermo (2008) son web sin paginación.
% Todas las claves están en core.bib (richards2015, richardsford2020,
% martin2017, cockburn2005, palermo2008): fuentes.bib no declara nada.
% Comillas: se usan comillas tipográficas literales, no \enquote, para no
% depender de csquotes (no está en la lista de paquetes de preamble.tex, #13).
% =============================================================================

% Grupo local: \emergencystretch evita líneas desbordadas sin afectar a otros capítulos.
\begingroup\emergencystretch=3em \hbadness=2000
\section*{Notas de investigación: Open Layers, Sinkhole Anti-pattern y atributos de calidad en N-Capas}

\section{Síntesis}

La arquitectura N-Capas organiza el sistema en capas horizontales en las que cada capa conoce
únicamente a la inmediatamente inferior. Una capa es \textbf{cerrada} cuando ninguna solicitud
puede saltársela —toda petición debe atravesarla para alcanzar la siguiente— y \textbf{abierta}
cuando sí puede evitarse, accediendo directamente a una capa más profunda
\parencite[3]{richards2015}. Las capas cerradas son las que materializan el \emph{aislamiento de
capas}: los cambios realizados en una capa no impactan a los componentes de las demás mientras
se respeten los contratos, de modo que cada capa es independiente y conoce poco o nada del
interior de las otras \parencite[3-4]{richards2015}.

Abrir una capa compra rendimiento a costa de ese aislamiento. El caso típico es el
\emph{fast-lane reader}: se declaran abiertas las capas de negocio y persistencia para que la
presentación lea directamente la base de datos en consultas simples de recuperación
\parencite[cap.~10]{richardsford2020}. El precio es un acoplamiento transversal —un cambio en
persistencia pasa a impactar a la vez a negocio y a presentación— y una arquitectura
descrita como frágil y cara de cambiar \parencite[4]{richards2015}. Por eso la decisión es
sobre todo de gobernanza: no documentar ni comunicar qué capas son abiertas y cuáles cerradas,
y por qué, produce arquitecturas acopladas y frágiles, difíciles de probar, mantener y desplegar
\parencite[5]{richards2015}.

El \emph{architecture sinkhole anti-pattern} aparece cuando las solicitudes atraviesan las capas
como simple \emph{pass-through}, sin que ninguna aporte lógica de negocio: el coste es
instanciación y procesamiento innecesarios, con impacto en memoria y rendimiento
\parencite[cap.~10]{richardsford2020}. Importa la dirección causal: las fuentes no presentan las
capas abiertas como causa del sinkhole, sino como reacción propuesta cuando el sinkhole domina
el tráfico, con el aviso explícito de que entonces será más difícil controlar el cambio por la
falta de aislamiento \parencite[7]{richards2015}. La \textbf{regla del 80/20} es el criterio de
diagnóstico: es normal que en torno al 20\,\% de las peticiones sean \emph{pass-through} y el
80\,\% porten lógica, y si la proporción se invierte eso indica que la arquitectura en capas no
es el estilo correcto para ese dominio \parencite[7]{richards2015}.

Las calificaciones bajas de la N-Capas en agilidad, facilidad de despliegue y escalabilidad no
son accidentales: proceden de su naturaleza monolítica y del acoplamiento entre componentes
\parencite[8-9]{richards2015}. Para requisitos de alta escalabilidad y agilidad, la Hexagonal y
la Limpia son estructuralmente más adecuadas porque invierten la dirección de las dependencias:
el núcleo declara los puertos y la infraestructura los implementa \parencite{cockburn2005}, con
dependencias que apuntan solo hacia adentro \parencite[203]{martin2017} y un dominio que no se
acopla más que a sí mismo \parencite{palermo2008}.

\section{Conceptos clave}

\begin{table}[htbp]
  \centering
  \footnotesize
  \renewcommand{\arraystretch}{1.2}
  \begin{tabular}{@{}p{0.22\linewidth}p{0.46\linewidth}p{0.24\linewidth}@{}}
    \toprule
    \textbf{Concepto} & \textbf{Definición} & \textbf{Fuente} \\
    \midrule
    \raggedright Capa cerrada
      & \raggedright Capa que no puede saltarse: toda solicitud debe atravesarla para llegar a
        la capa inferior
      & \raggedright \textcite[3]{richards2015}; \textcite[cap.~10]{richardsford2020}
        \tabularnewline
    \raggedright Capa abierta
      & \raggedright Capa que una solicitud puede evitar, accediendo directamente a una capa más
        profunda
      & \raggedright \textcite[3]{richards2015}; \textcite[cap.~10]{richardsford2020}
        \tabularnewline
    \raggedright Aislamiento de capas
      & \raggedright Los cambios en una capa no impactan a las demás mientras se mantengan los
        contratos entre capas
      & \raggedright \textcite[3-4]{richards2015}; \textcite[cap.~10]{richardsford2020}
        \tabularnewline
    \raggedright Gobernanza de las capas abiertas
      & \raggedright No documentar qué capas son abiertas o cerradas produce arquitecturas
        acopladas y frágiles, difíciles de probar, mantener y desplegar
      & \raggedright \textcite[5]{richards2015} \tabularnewline
    \raggedright \emph{Fast-lane reader}
      & \raggedright Declarar abiertas negocio y persistencia para que la presentación lea
        directamente la base de datos en consultas simples de alto volumen
      & \raggedright \textcite[cap.~10]{richardsford2020} \tabularnewline
    \raggedright \emph{Architecture sinkhole anti-pattern}
      & \raggedright Solicitudes que atraviesan las capas como \emph{pass-through}, sin lógica
        en ninguna, con coste de instanciación y procesamiento
      & \raggedright \textcite[7]{richards2015}; \textcite[cap.~10]{richardsford2020}
        \tabularnewline
    \raggedright Regla del 80/20
      & \raggedright Si en torno al 80\,\% de las peticiones son \emph{pass-through}, el estilo
        en capas no es el adecuado para ese dominio
      & \raggedright \textcite[7]{richards2015}; \textcite[cap.~10]{richardsford2020}
        \tabularnewline
    \raggedright Partición técnica frente a partición por dominio
      & \raggedright En N-Capas el dominio queda disperso entre capas técnicas, por lo que un
        enfoque dirigido por el dominio encaja mal
      & \raggedright \textcite[cap.~10]{richardsford2020} \tabularnewline
    \raggedright Calificaciones de la N-Capas
      & \raggedright Agilidad y facilidad de despliegue bajas por la naturaleza monolítica y el
        acoplamiento; escalabilidad baja por la misma razón
      & \raggedright \textcite[8-9]{richards2015}; \textcite[cap.~10]{richardsford2020}
        \tabularnewline
    \raggedright Regla de Dependencia
      & \raggedright Las dependencias del código fuente apuntan solo hacia adentro, hacia las
        políticas de mayor nivel
      & \raggedright \textcite[203]{martin2017} \tabularnewline
    \raggedright Testabilidad del núcleo
      & \raggedright Las reglas de negocio se prueban sin interfaz de usuario, base de datos ni
        servidor
      & \raggedright \textcite[202]{martin2017} \tabularnewline
    \raggedright Puertos y adaptadores; dominio al centro
      & \raggedright El núcleo declara los puertos y la infraestructura los implementa; el
        Modelo de Dominio no se acopla más que a sí mismo
      & \raggedright \textcite{cockburn2005}; \textcite{palermo2008} \tabularnewline
    \bottomrule
  \end{tabular}
  \caption{Conceptos clave de la investigación INV-P6.}
  \label{tab:inv-p6-conceptos}
\end{table}

\section{Puntos de conexión con las preguntas}

\begin{itemize}
  \item \textbf{P6, primer apartado:} capa cerrada frente a capa abierta y aislamiento de capas
        dan el mecanismo exacto por el que abrir una capa degrada la mantenibilidad
        \parencite[3-5]{richards2015}.
  \item \textbf{P6, segundo apartado:} el sinkhole y la regla del 80/20 permiten mostrar que la
        apertura de capas es reacción al antipatrón, no su causa, y que la regla es un
        diagnóstico de encaje del estilo, no una perilla de optimización
        \parencite[7]{richards2015}.
  \item \textbf{P6, tercer apartado:} las calificaciones bajas \parencite[8-9]{richards2015} y
        la partición técnica \parencite[cap.~10]{richardsford2020} se contraponen a la inversión
        de dependencias de la Hexagonal y la Limpia \parencite[203]{martin2017}.
  \item \textbf{P1 y P2:} el mismo contraste N-Capas frente a Hexagonal/Onion sostiene el
        análisis de \emph{trade-offs} de P1 y la comparación de simetría de P2.
\end{itemize}

\section*{Notas de diseño}

\begin{itemize}
  \item \textbf{Dirección causal del sinkhole.} El enunciado de P6 sugiere que las capas
        abiertas llevan al sinkhole; las fuentes dicen lo inverso: son la reacción propuesta
        cuando el sinkhole domina \parencite[7]{richards2015}. La investigación documenta el
        bucle completo —capas vacías, sinkhole, apertura sancionada del rodeo, vaciado
        acelerado— para que la respuesta pueda ser fiel a las fuentes sin contradecir el
        enunciado.
  \item \textbf{Paginación de Richards y Ford (2020).} El PDF de \texttt{Bibliografias/} no
        conserva la paginación impresa al extraer el texto, así que ese libro se cita por
        capítulo y sección. Las páginas de \textcite{richards2015} se verificaron contra la
        numeración impresa del capítulo~1, que en el PDF local aparece desplazada por las
        páginas preliminares.
  \item \textbf{Doble cita de Richards.} \textcite{richards2015} aporta las calificaciones del
        patrón y \textcite{richardsford2020} el tratamiento ampliado del mismo autor principal
        (capas abiertas, \emph{fast-lane reader}, partición técnica). Se citan juntos donde
        coinciden y por separado donde cada uno es más explícito.
  \item \textbf{De la Torre Llorente y otros (2010).} Figura en \texttt{fuentes.md} como
        contexto del seminario sobre N-Capas orientada al dominio, pero no aporta ninguna
        afirmación de estas notas, así que no se cita.
  \item \textbf{DDD sin Evans.} El enunciado menciona principios de DDD. Se sostienen con la
        partición por dominio \parencite[cap.~10]{richardsford2020} y el dominio al centro
        \parencite{palermo2008}, sin recurrir a Evans, que no figura en la bibliografía
        consultada para esta investigación.
\end{itemize}
\endgroup
.
% Estructura del issue #12: \section*{Notas de investigación: …} + Síntesis +
% Conceptos clave (tabla booktabs) + Puntos de conexión + Notas de diseño.
% Fuente de trabajo: notas.md (misma carpeta).
% Citas: páginas impresas de Richards (2015) y Martin (2017); Richards y Ford
% (2020) no conserva paginación al extraer el texto del PDF y se cita por
% capítulo; Cockburn (2005) y Palermo (2008) son web sin paginación.
% Todas las claves están en core.bib (richards2015, richardsford2020,
% martin2017, cockburn2005, palermo2008): fuentes.bib no declara nada.
% Comillas: se usan comillas tipográficas literales, no \enquote, para no
% depender de csquotes (no está en la lista de paquetes de preamble.tex, #13).
% =============================================================================

% Grupo local: \emergencystretch evita líneas desbordadas sin afectar a otros capítulos.
\begingroup\emergencystretch=3em \hbadness=2000
\section*{Notas de investigación: Open Layers, Sinkhole Anti-pattern y atributos de calidad en N-Capas}

\section{Síntesis}

La arquitectura N-Capas organiza el sistema en capas horizontales en las que cada capa conoce
únicamente a la inmediatamente inferior. Una capa es \textbf{cerrada} cuando ninguna solicitud
puede saltársela —toda petición debe atravesarla para alcanzar la siguiente— y \textbf{abierta}
cuando sí puede evitarse, accediendo directamente a una capa más profunda
\parencite[3]{richards2015}. Las capas cerradas son las que materializan el \emph{aislamiento de
capas}: los cambios realizados en una capa no impactan a los componentes de las demás mientras
se respeten los contratos, de modo que cada capa es independiente y conoce poco o nada del
interior de las otras \parencite[3-4]{richards2015}.

Abrir una capa compra rendimiento a costa de ese aislamiento. El caso típico es el
\emph{fast-lane reader}: se declaran abiertas las capas de negocio y persistencia para que la
presentación lea directamente la base de datos en consultas simples de recuperación
\parencite[cap.~10]{richardsford2020}. El precio es un acoplamiento transversal —un cambio en
persistencia pasa a impactar a la vez a negocio y a presentación— y una arquitectura
descrita como frágil y cara de cambiar \parencite[4]{richards2015}. Por eso la decisión es
sobre todo de gobernanza: no documentar ni comunicar qué capas son abiertas y cuáles cerradas,
y por qué, produce arquitecturas acopladas y frágiles, difíciles de probar, mantener y desplegar
\parencite[5]{richards2015}.

El \emph{architecture sinkhole anti-pattern} aparece cuando las solicitudes atraviesan las capas
como simple \emph{pass-through}, sin que ninguna aporte lógica de negocio: el coste es
instanciación y procesamiento innecesarios, con impacto en memoria y rendimiento
\parencite[cap.~10]{richardsford2020}. Importa la dirección causal: las fuentes no presentan las
capas abiertas como causa del sinkhole, sino como reacción propuesta cuando el sinkhole domina
el tráfico, con el aviso explícito de que entonces será más difícil controlar el cambio por la
falta de aislamiento \parencite[7]{richards2015}. La \textbf{regla del 80/20} es el criterio de
diagnóstico: es normal que en torno al 20\,\% de las peticiones sean \emph{pass-through} y el
80\,\% porten lógica, y si la proporción se invierte eso indica que la arquitectura en capas no
es el estilo correcto para ese dominio \parencite[7]{richards2015}.

Las calificaciones bajas de la N-Capas en agilidad, facilidad de despliegue y escalabilidad no
son accidentales: proceden de su naturaleza monolítica y del acoplamiento entre componentes
\parencite[8-9]{richards2015}. Para requisitos de alta escalabilidad y agilidad, la Hexagonal y
la Limpia son estructuralmente más adecuadas porque invierten la dirección de las dependencias:
el núcleo declara los puertos y la infraestructura los implementa \parencite{cockburn2005}, con
dependencias que apuntan solo hacia adentro \parencite[203]{martin2017} y un dominio que no se
acopla más que a sí mismo \parencite{palermo2008}.

\section{Conceptos clave}

\begin{table}[htbp]
  \centering
  \footnotesize
  \renewcommand{\arraystretch}{1.2}
  \begin{tabular}{@{}p{0.22\linewidth}p{0.46\linewidth}p{0.24\linewidth}@{}}
    \toprule
    \textbf{Concepto} & \textbf{Definición} & \textbf{Fuente} \\
    \midrule
    \raggedright Capa cerrada
      & \raggedright Capa que no puede saltarse: toda solicitud debe atravesarla para llegar a
        la capa inferior
      & \raggedright \textcite[3]{richards2015}; \textcite[cap.~10]{richardsford2020}
        \tabularnewline
    \raggedright Capa abierta
      & \raggedright Capa que una solicitud puede evitar, accediendo directamente a una capa más
        profunda
      & \raggedright \textcite[3]{richards2015}; \textcite[cap.~10]{richardsford2020}
        \tabularnewline
    \raggedright Aislamiento de capas
      & \raggedright Los cambios en una capa no impactan a las demás mientras se mantengan los
        contratos entre capas
      & \raggedright \textcite[3-4]{richards2015}; \textcite[cap.~10]{richardsford2020}
        \tabularnewline
    \raggedright Gobernanza de las capas abiertas
      & \raggedright No documentar qué capas son abiertas o cerradas produce arquitecturas
        acopladas y frágiles, difíciles de probar, mantener y desplegar
      & \raggedright \textcite[5]{richards2015} \tabularnewline
    \raggedright \emph{Fast-lane reader}
      & \raggedright Declarar abiertas negocio y persistencia para que la presentación lea
        directamente la base de datos en consultas simples de alto volumen
      & \raggedright \textcite[cap.~10]{richardsford2020} \tabularnewline
    \raggedright \emph{Architecture sinkhole anti-pattern}
      & \raggedright Solicitudes que atraviesan las capas como \emph{pass-through}, sin lógica
        en ninguna, con coste de instanciación y procesamiento
      & \raggedright \textcite[7]{richards2015}; \textcite[cap.~10]{richardsford2020}
        \tabularnewline
    \raggedright Regla del 80/20
      & \raggedright Si en torno al 80\,\% de las peticiones son \emph{pass-through}, el estilo
        en capas no es el adecuado para ese dominio
      & \raggedright \textcite[7]{richards2015}; \textcite[cap.~10]{richardsford2020}
        \tabularnewline
    \raggedright Partición técnica frente a partición por dominio
      & \raggedright En N-Capas el dominio queda disperso entre capas técnicas, por lo que un
        enfoque dirigido por el dominio encaja mal
      & \raggedright \textcite[cap.~10]{richardsford2020} \tabularnewline
    \raggedright Calificaciones de la N-Capas
      & \raggedright Agilidad y facilidad de despliegue bajas por la naturaleza monolítica y el
        acoplamiento; escalabilidad baja por la misma razón
      & \raggedright \textcite[8-9]{richards2015}; \textcite[cap.~10]{richardsford2020}
        \tabularnewline
    \raggedright Regla de Dependencia
      & \raggedright Las dependencias del código fuente apuntan solo hacia adentro, hacia las
        políticas de mayor nivel
      & \raggedright \textcite[203]{martin2017} \tabularnewline
    \raggedright Testabilidad del núcleo
      & \raggedright Las reglas de negocio se prueban sin interfaz de usuario, base de datos ni
        servidor
      & \raggedright \textcite[202]{martin2017} \tabularnewline
    \raggedright Puertos y adaptadores; dominio al centro
      & \raggedright El núcleo declara los puertos y la infraestructura los implementa; el
        Modelo de Dominio no se acopla más que a sí mismo
      & \raggedright \textcite{cockburn2005}; \textcite{palermo2008} \tabularnewline
    \bottomrule
  \end{tabular}
  \caption{Conceptos clave de la investigación INV-P6.}
  \label{tab:inv-p6-conceptos}
\end{table}

\section{Puntos de conexión con las preguntas}

\begin{itemize}
  \item \textbf{P6, primer apartado:} capa cerrada frente a capa abierta y aislamiento de capas
        dan el mecanismo exacto por el que abrir una capa degrada la mantenibilidad
        \parencite[3-5]{richards2015}.
  \item \textbf{P6, segundo apartado:} el sinkhole y la regla del 80/20 permiten mostrar que la
        apertura de capas es reacción al antipatrón, no su causa, y que la regla es un
        diagnóstico de encaje del estilo, no una perilla de optimización
        \parencite[7]{richards2015}.
  \item \textbf{P6, tercer apartado:} las calificaciones bajas \parencite[8-9]{richards2015} y
        la partición técnica \parencite[cap.~10]{richardsford2020} se contraponen a la inversión
        de dependencias de la Hexagonal y la Limpia \parencite[203]{martin2017}.
  \item \textbf{P1 y P2:} el mismo contraste N-Capas frente a Hexagonal/Onion sostiene el
        análisis de \emph{trade-offs} de P1 y la comparación de simetría de P2.
\end{itemize}

\section*{Notas de diseño}

\begin{itemize}
  \item \textbf{Dirección causal del sinkhole.} El enunciado de P6 sugiere que las capas
        abiertas llevan al sinkhole; las fuentes dicen lo inverso: son la reacción propuesta
        cuando el sinkhole domina \parencite[7]{richards2015}. La investigación documenta el
        bucle completo —capas vacías, sinkhole, apertura sancionada del rodeo, vaciado
        acelerado— para que la respuesta pueda ser fiel a las fuentes sin contradecir el
        enunciado.
  \item \textbf{Paginación de Richards y Ford (2020).} El PDF de \texttt{Bibliografias/} no
        conserva la paginación impresa al extraer el texto, así que ese libro se cita por
        capítulo y sección. Las páginas de \textcite{richards2015} se verificaron contra la
        numeración impresa del capítulo~1, que en el PDF local aparece desplazada por las
        páginas preliminares.
  \item \textbf{Doble cita de Richards.} \textcite{richards2015} aporta las calificaciones del
        patrón y \textcite{richardsford2020} el tratamiento ampliado del mismo autor principal
        (capas abiertas, \emph{fast-lane reader}, partición técnica). Se citan juntos donde
        coinciden y por separado donde cada uno es más explícito.
  \item \textbf{De la Torre Llorente y otros (2010).} Figura en \texttt{fuentes.md} como
        contexto del seminario sobre N-Capas orientada al dominio, pero no aporta ninguna
        afirmación de estas notas, así que no se cita.
  \item \textbf{DDD sin Evans.} El enunciado menciona principios de DDD. Se sostienen con la
        partición por dominio \parencite[cap.~10]{richardsford2020} y el dominio al centro
        \parencite{palermo2008}, sin recurrir a Evans, que no figura en la bibliografía
        consultada para esta investigación.
\end{itemize}
\endgroup
